% Appendix: project record, 24 Oct to 14 Nov 2026 (writer id midnov). Draft 27 Nov 2026.
% Every number and dated fact is traced in app_midnov_NOTES.md (file, line, quote). No paths or host names here.
\chapter{From the SSL series to the benchmark (24 Oct to 14 Nov)}\label{app:midnov}

This appendix records, in date order, the work between the end of the semi-supervised series (the AL1 run of 23 Oct) and the
start of the benchmark build on 12 Nov, up to E1 opening \#7 on 14 Nov. Chapter~\ref{ch:legacy} summarises this period, and
Table~\ref{tab:T21} (one row per arm) and Table~\ref{tab:T17} (one row per E1 opening) list its outcomes. Here each item gets the
fuller record: what was asked, what was run, the result, the verdict under the rule in force at the time, and what it led to,
including stopped, failed and invalid runs. Nothing here changes a row of those tables. A desk study of thermal imaging (4 Nov) ran
in a separate team and is not recorded here.

\section{Conditions of every number in this appendix}\label{sec:midnov-conditions}

\paragraph{Data and labels.} All runs use the legacy split (478 train, 54 validation, 50 test cells) and train on label set v1
(Section~\ref{sec:data-labels}) unless stated. \emph{test50} is the 50 legacy test cells; \emph{clean 49} is test50 without the one
very noisy cell. test50 overlaps the training pool (Section~\ref{sec:data-split}) and was scored at every rung and every opening of
this period, so \textbf{every test50 value below is in-distribution, with known leakage, and context only}. E1 is the 106 cells of
Section~\ref{sec:data-sets}; every E1 value below is a repeated look at the same 106 cells.

\paragraph{Metrics.} Crack IoU and inactive IoU are per-class IoU at argmax from one confusion matrix pooled over all pixels of the
scored cell set (the standard convention; Section~\ref{sec:bg_metrics}). relF1@2, relP@2 and relR@2 are the 2-px tolerance metrics
of the same section, pooled over the set. mIoU\_rob (Track M) is the mean IoU over the 13 classes with ground truth in at least 3
test50 cells, background included. Crack AP is the binned pixel average precision of the crack softmax probability, pooled over
the set. All values are on a 0 to 1 scale.

\paragraph{Rules.} \emph{Rule 1} (house rule, locked before each run): a difference between two arms of seed means is BETTER or
WORSE if its size exceeds $\max(0.01,\, 2\sqrt{(s_A^2 + s_B^2)/3})$, with $s_A, s_B$ the seed standard deviations; otherwise it is
n.d.\ (not detectable). An arm passes if it is BETTER on at least one primary and WORSE on none, read on two verdict sets: all 50 and
clean 49. Ensemble contrasts use a paired image bootstrap on test50 or a 71-cluster bootstrap on E1 (2,000 resamples) and pass when
the 95\% interval lies wholly above 0. \emph{Guards} are pre-registered FAIL rules: a tripped guard fails the arm whatever else
passes. Below, the \emph{verdict} is always the rule fixed before the run; values printed beside it decide nothing; interpretations
are marked ``our reading''.

\paragraph{Machines.} The \emph{training machine} (eight shared GPUs) and the \emph{second machine} (one GPU, shared with another
user's process, for which at least 2\,GB had to stay free) ran different arms. Seeds from the two machines are never pooled; the 100\%-label
baseline was therefore retrained on the second machine (SUP100-g) before any comparison there.

\paragraph{E1 policy.} Until 31 Oct E1 was read many times as a descriptive set (P-series, Track M, the 27 Oct team, the
inference checks of 30 Oct, the protocol\_100 openings). From about 15:00 on 31 Oct, by the user's decision, E1 became a second
test set: a protocol could pre-register E1 co-primaries, read once at its single opening, never for tuning or selection. The models
of FINAL and of protocol\_100 (seeds 0 to 4) are therefore not blind on E1.

\section{The E-SCDD download and audit (12 Oct; addition of 30 Oct)}\label{sec:midnov-escdd}

The E-SCDD extension \citep{hanifi2025escdd} was downloaded in full on 12 Oct and audited without training a model. Track M's R4 rung later trained on 107 of its
real cells, and its scuff labels drove the scuff work of this period. Chapter~\ref{ch:data} gives the label sets that resulted.
\begin{itemize}
  \item \textbf{What was refined.} BenchmarkEL has three releases; the project uses 20211104 (582 cells). The maintainers'
    20221008 release (695 cells) corrected many crack labels (dotted labels joined, lines widened). E-SCDD repackages 20221008. On
    the same 265 cells, crack pixels rise from 327,754 to 396,572, pooled width from 2.17 to 2.52 px and fragments fall from 5,391
    to 1,288. The ``revised labels'' that moved crack IoU by $+0.05$ to $+0.08$ in the earlier work are these 20221008
    corrections, obtained through E-SCDD.
  \item \textbf{What E-SCDD itself changed.} All 693 shared images are pixel-identical to 20221008; 627 masks are identical and the
    other 66 differ only between cell-gap classes (62 masks from sp multi to the new class sp multi halfcut, 552,426 px; 4 swaps).
    Crack and inactive pixels changed: 0.
  \item \textbf{What was not as described.} The release holds 903 labelled files (693 + 117 new real + 93 synthetic) against the
    paper's 1,000 (695 + 110 + 195), and 0 of the unlabelled synthetic images the paper reports. The 93 synthetic images come from
    4 base cells, all in E-SCDD's train split, and hold 11.7\% of its crack pixels. The 117 new real cells (114 CSIR, 3 ARTS) carry
    less consistent crack labels (3 cells with 128 to 380 fragments); 100 of them are in our unlabelled pool.
  \item \textbf{Splits.} E-SCDD's split is not ours: 48 of its 82 test cells are in our train split and 39 of our 50 test cells in
    its train split, so neither test set can score models trained on the other split.
  \item \textbf{Scuff (added 30 Oct).} A CPU read of the raw E-SCDD masks found 12 train cells with scuff (class 25, 3,246 px).
    Our v1 build had mapped those pixels back to their old labels: 1,207 background, 1,390 material and 649 crack in 6 cells. The
    one ARTS crack cell of the 5\% label set (47 px) is entirely scuff. In test50, 236 scuff px in 2 cells are crack. These facts
    led to the scuff class arm (B3, Section~\ref{sec:midnov-thin}) and later to label set v1.1.
\end{itemize}

\section{Two analysis teams, the goal, Track D and Track M (24 to 27 Oct)}\label{sec:midnov-tracks}

\paragraph{Why-SOTA team (24 Oct, no GPU).} After AL1, two analysts and a critic asked why 2025 SSL methods fail on EL crack. From
CPU recomputes of existing outputs: SSL was not doing nothing on crack (BsM5 beat supervised by $+0.029$ crack IoU at 5\% labels and
still passed with a free logit bias (logit adjustment) on both sides, $+0.037$), but the gain was a third of what 24 random labelled cells buy
($+0.092$) and was source-bound. About 38\% of pool crack was a true miss, more than 2\,px from any prediction. The critic ranked
eight directions; the first was to stop new SSL methods for crack. The user then set the goal: ``highest scores, a model which
captures almost all of the cracks, inactive and other important defects balancing the miou for other classes as well'', and asked
the team to act on the critic's decisions. That started Track D (diagnosis) and Track M (best model), both on the training machine.

\paragraph{Track D (24 Oct).} Three diagnostics, each with a rule fixed before the run (Table~\ref{tab:midnov-trackd}).
D1 was a linear probe on frozen encoder features, trained on the 23 labelled cells of the 5\% set; D2 re-scored existing models with
crack within 2\,px of inactive ignored; D3 added 12 cells to the 5\% base in three ways, 3 draws $\times$ 3 seeds per arm (27 runs),
with the cells chosen using their labels (oracle-selected). D3's 3-seed design broke the house rule of at least 5 seeds; the user
accepted it with the critic's decisions.

\begin{table}[htbp]
\centering
\footnotesize
\caption[Track D diagnostics (24 Oct)]{Track D diagnostics (24 Oct; training machine; 5\% label base; test50, in-distribution,
known leakage, context only). BMRK crack AP: crack AP over the BMRK test cells.}
\label{tab:midnov-trackd}
\begin{tabular}{@{}p{0.17\linewidth}p{0.33\linewidth}p{0.42\linewidth}@{}}
\toprule
Item & Rule (fixed before) & Result and verdict \\
\midrule
D1 frozen probe & pass if BMRK crack AP beats the 5\% model with logit bias (0.165) by 0.05, bar 0.215 & ResNet-50 0.268 (pass on test50 and pool); ConvNeXt-T 0.278 / 0.271 (pass on test50 only); ResNet-34 0.185, DINOv2-S 0.122, the 5\% model's own encoder 0.171 (fail). Beside: at the 5\% model's precision the probes recall 0.020 to 0.023 (ResNet-50) and 0.002 (ConvNeXt-T) of BMRK crack against its 0.21 \\
D2 inactive edge & kill if a banded SSL gap passes its threshold with the opposite sign & band holds 11.6\% of test crack; banding raises 5 to 10\% arms by $+0.018$ to $+0.027$ crack IoU; no sign flip: kill not triggered \\
D3 labelled set & draw-level rule 1 over 3 draw means; P1 Ts $-$ Tc, P2 Tc $-$ R & P2 crack cells $-$ random: crack IoU $+0.0127$ (thr 0.0118) BETTER; P1 weak-source crack cells $-$ crack cells: $+0.0151$ (thr 0.0212) n.d. Beside: Tc $-$ R relF1@2 $+0.030$ $[+0.020, +0.042]$ in all four cell sets \\
\bottomrule
\end{tabular}
\end{table}

D1's rule closed masked-image pretraining for the time being. D3's gate for SSL on the better-covered set (Ts $-$ Tc BETTER, or Tc
$-$ R BETTER with BMRK recall $\geq +0.05$) was not met (BMRK recall $+0.023$), so that SSL arm never ran. Our reading: at this budget
targeted labels beat random labels and SSL for crack, but only with oracle selection, and crack-pixel volume explains the Tc $-$ R
gap.

\paragraph{Track M (24 to 25 Oct).} The baseline SUP100 (all 478 labels, U-Net++ with ImageNet ResNet-34, cross-entropy plus Dice on
crack and inactive, label-safe augmentation, 512\,px, 20,000 steps, 5 seeds) was first re-scored (M0): crack IoU $0.435 \pm 0.078$
on all 50 and $0.474 \pm 0.002$ on clean 49; one seed failed on the noisy cell (0.297), so every later contrast was read with and
without that cell. Then one change per rung was tried, each entering the final recipe only if it passed rule 1 on validation;
test50 was scored per rung but never used to choose (Table~\ref{tab:midnov-trackm}).

\begin{table}[htbp]
\centering
\footnotesize
\caption[Track M rungs at 100\% labels (24 to 25 Oct)]{Track M rungs at 100\% labels (training machine; 5 v 5 seeds; deltas on
clean 49 against SUP100, threshold in parentheses; test50, in-distribution, known leakage, context only). Track M test openings
1 to 7 in the order R1, R2, R4, R2c, R3 Part B, R5, FINAL.}
\label{tab:midnov-trackm}
\begin{tabular}{@{}p{0.27\linewidth}p{0.45\linewidth}p{0.20\linewidth}@{}}
\toprule
Rung & Result & Verdict \\
\midrule
R1 TTA, ensemble, crack bias & TTA $-$ plain $+0.0104$ (0.0100); 5-seed ensemble $-$ mean single $+0.0144$ $[+0.0118, +0.0167]$; validation chose bias $+0.75$ & R1a PASS (clean 49; all 50 n.d.); R1b PASS; validation set TTA off for the ensemble \\
R2 soft Dice on 12 classes & crack IoU $-0.0106$ (0.0100), relF1@2 $-0.0130$ (0.0100); gridline 0.283 to 0.381 & WORSE; excluded \\
R2c Dice on crack, inactive, gridline & n.d.\ on all four primaries; gridline 0.283 to 0.357 & excluded \\
R3 U-Net++ ResNet-50 (pilot picked over EfficientNet-B4 and ConvNeXt-T on validation) & inactive IoU $-0.0251$ (0.0227); crack $-0.0041$ & WORSE; excluded \\
R4 +107 E-SCDD cells (104 CSIR) & crack $+0.0010$ $[-0.0026, +0.0062]$, inactive $+0.0085$ $[-0.0042, +0.0277]$ & n.d.; excluded \\
R5 2$\times$ internal resolution & crack $-0.0046$ (0.0100) at 4.0$\times$ compute & n.d.; excluded \\
\bottomrule
\end{tabular}
\end{table}

R2's crack loss came with a gradient clip rate of 0.969 against 0.253 (mean pre-clip norm 6.74 against 1.80), so it is attributed to Dice
on 12 classes at the house clip, not to Dice in general. No training change passed, so by the locked protocol the final recipe was
SUP100 itself, and FINAL is its 5-seed mean-softmax ensemble with crack logit $+0.75$ (chosen as the maximum of validation relF1@2,
since no bias reached relR@2 $\geq$ 0.90 at relP@2 $\geq$ 0.90). \textbf{Verdict (opening 7, 25 Oct): PASS} against the mean single
seed on crack IoU and relF1@2 at both biases: crack IoU on clean 49 $+0.018$ $[+0.013, +0.025]$ (0.474 to 0.492). Beside it: on the
noisy cell bias $+0.75$ calls 10,867 px crack (bias 0 calls 2), so bias 0 is the robust operating point; relR@2 / relP@2 reached
0.866 / 0.930, so ``almost all cracks'' at 0.90 / 0.90 was not reached. ResNet-50 and 2$\times$ resolution both raised E1 crack
recall (0.178 to 0.214 and 0.204), descriptive only. A finding of 30 Oct (Section~\ref{sec:midnov-erasure}) qualifies every
``no change helps'' of this track: the optimizer was never varied.

\paragraph{Unlabelled-data team (27 Oct, no GPU).} The user asked why about 12k unlabelled cells had not helped. The team's answer:
only label propagation had been tested, never feature learning from the pool, and it had worked only with few labels. On the clean
cells a crack gain of $+0.02$ or more was excluded at 20, 50 and 100\% labels; the old ``$+0.028$ at 100\%'' had been one supervised
seed failing on the noisy cell. Of FINAL's crack IoU loss, 70\% lay within 2\,px of the label. On E1, 20 half-cut cells with a
scuff streak held 75\% of the crack pixels, and FINAL found 0.126 of their crack against 0.552 on half-cut cells without scuff: the
model painted the streak as crack and missed the cracks branching off it. The team advised keeping test50, adding the clean set as
co-primary, keeping E1 sealed with two readouts (scuff plus half-cut, and the rest), and ranked five next steps; item 3 was one test
of the user's own intuition: masked-image pretraining (MIM) of the encoder, fine-tuned at 5\% labels.

\section{Feature learning from the pool at 5\% labels (28 to 30 Oct)}\label{sec:midnov-features}

On 28 Oct the user asked whether learning features from the unlabelled cells, without labels, improves the model trained on all
labels; the answer was that it had never been run. A new team ran three methods on the second machine, against a fresh 5\%
supervised baseline (A5c-f1, 5 seeds; test50 crack IoU 0.300 against 0.304 for the original on the training machine). All three
were screened at 5\% labels, a choice copied from the 27 Oct plan without checking it against the user's question, which concerned
100\% labels (Section~\ref{sec:midnov-erasure}).
\begin{itemize}
  \item \textbf{MIM pretraining.} SparK-style sparse masked-image modelling \citep{tian2023spark} of the ImageNet ResNet-34
    encoder on a pool of 11,690 cells (every module of a validation, test or E1 cell excluded, plus near-duplicates), mask ratio 0.6,
    batch 32, 36,500 steps (29 Oct 13:08 to 18:34). Monitor loss 0.914 to 0.208.
  \item \textbf{MIM5} (MIM encoder, house fine-tune, 5 v 5). \textbf{Verdict: n.d.} on all three primaries in both sets, every
    estimate negative (crack IoU 0.3002 to 0.2886, $-0.0116$, thr 0.0186). Closed per protocol.
  \item \textbf{AM} (MIM reconstruction error as a crack map, validation only). \textbf{KILL}: map AP 0.0166; no threshold reached
    relaxed precision 0.5, so no missed component was recovered.
  \item \textbf{DDeP5} (decoder denoising pretraining \citep{brempong2022ddep}, 40,000 steps on the training machine, then 5 v 5).
    \textbf{Verdict: n.d.}, small positive estimates (crack IoU $+0.0070$, thr 0.0100), mostly a crack threshold shift (recall 0.365
    to 0.384, precision 0.628 to 0.606). Closed per protocol.
  \item \textbf{CP5} (labelled crack components pasted onto pool cells, 20-component library). Its decoy gate passed at 2,000 steps;
    the arm was \textbf{stopped unscored} at 10:50 on 30 Oct when the 5\% choice was recognised as wrong (two seeds finished, two
    stopped at 11,800 and 12,350 steps, one cancelled). No test50 number was opened. Its only ARTS paste was the scuff cell of
    Section~\ref{sec:midnov-escdd}.
\end{itemize}

\section{The encoder audit and the encoder-preserving fine-tune (30 Oct)}\label{sec:midnov-erasure}

On 30 Oct a direct read of the weights (Table~\ref{tab:midnov-cos}) showed that the house fine-tune (Adam, learning rate
$10^{-3}$, coupled weight decay $10^{-4}$, 20,000 steps, never varied in the P-series or Track M) leaves encoder layer3 and layer4
orthogonal to their initial weights, for ImageNet and MIM initialisations, at 5\% and 100\% labels. In a 2,000-step smoke of SUP100
the cosines of layer3 and layer4 to ImageNet were already 0.070 and 0.003, while their norms were still 0.33 and 0.21 of initial: the
direction is lost during warm-up and peak learning rate, before the norms shrink. MIM had changed mainly layer2 to layer4, so MIM5
tested only the transfer of conv1 and layer1. Every earlier ``initialisation, pretraining or backbone does not help'' result (P15,
MIM5, DDeP5, Track M R3) was obtained under this optimizer; erasure was measured on ResNet-34 U-Net++ runs only.

\begin{table}[htbp]
\centering
\footnotesize
\caption[Encoder audit: cosine of final to initial weights per stage (30 Oct)]{Encoder audit, 30 Oct: per-tensor cosine of encoder
convolution weights, averaged per ResNet-34 stage, at step 20,000 (first column: the MIM encoder after pretraining; training only, no test number involved).}
\label{tab:midnov-cos}
\begin{adjustbox}{max width=\linewidth}
\begin{tabular}{@{}lccccc@{}}
\toprule
Stage & MIM vs ImageNet & A5c-f1 final vs ImageNet & MIM5 final vs MIM & SUP100 s0 final vs ImageNet & norm final / initial (SUP100) \\
\midrule
conv1 & 0.996 & 0.980 & 0.976 & 0.985 & 0.86 \\
layer1 & 0.859 & 0.379 & 0.327 & 0.396 & 0.75 \\
layer2 & 0.761 & 0.135 & 0.139 & 0.150 & 0.60 \\
layer3 & 0.618 & 0.010 & 0.005 & 0.018 & 0.09 \\
layer4 & 0.533 & $-0.001$ & $-0.000$ & 0.001 & 0.06 \\
\bottomrule
\end{tabular}
\end{adjustbox}
\end{table}

\paragraph{Why 5\% was the wrong test.} The user, about 10:45 on 30 Oct: ``why 5\%? I asked about adding the pool to 100\% labels.''
An analysis that day found that a null at 5\% carries to 100\% only if the method does not consume labels and the 5\% estimate is
tight on crack: copy-paste meets neither (its library grows from about 20 to about 560 components), inactive never meets the second
at 5\% (2 training cells with inactive), and 100\% is better powered (minimum detectable crack IoU change 0.012 against 0.019 at 5\%).
Pool ideas were thereafter tested on top of 100\% labels.

\paragraph{PF.} The encoder-preserving fine-tune freezes the encoder for steps 0 to 999 (batch norm in train mode), then ramps the
encoder learning rate from 0 to 0.1 times the base over 1,000 steps, with encoder weight decay 0; decoder and head are trained as in
the house recipe. Its gate before launch required layer3 and layer4 cosine $\geq$ 0.8 at 4,000 steps: SUP100-PF 0.975 / 0.985,
MIM100-PF 0.986 / 0.990, pass.

\section{protocol\_100 and the label work (30 to 31 Oct)}\label{sec:midnov-p100}

protocol\_100 answered the user's question at 100\% labels: 4 arms $\times$ 5 seeds on the second machine, all 478 labels, 20,000
steps. SUP100-g (fresh house baseline on that machine), SUP100-PF, MIM100-PF (PF from the MIM encoder) and CP100 (house optimizer
plus a crack copy-paste stream from a frozen library of 554 components). Gates 1 to 4 (bit-identity checks of the new switches when off and of validation off)
passed; the copy-paste decoy gate at 2,000 steps was flagged (2.45 times the baseline's false crack on decoy spots against a 2$\times$
line), and an amendment made before any CP100 run (gate 6b) re-measured it at 4,000 steps, where it fell (0.0076 to 0.0039), so CP100
ran. The fine-tunes ran from 12:05 on 30 Oct to 00:56 on 31 Oct, CP100 from 00:21 to 06:52; a reboot of the machine at 16:28 stopped
three runs, which resumed from checkpoints with loss and learning-rate lines identical to those before the reboot. The two test
openings were scored at 00:57 to 01:33 and 06:53 to 07:07 on 31 Oct (Table~\ref{tab:midnov-p100}).

\begin{table}[htbp]
\centering
\footnotesize
\caption[protocol\_100 openings 1 and 2 (31 Oct)]{protocol\_100 openings 1 and 2 (second machine; 5 v 5; rule 1; delta (threshold)
on all 50 / clean 49; test50, in-distribution, known leakage, context only).}
\label{tab:midnov-p100}
\begin{adjustbox}{max width=\linewidth}
\begin{tabular}{@{}p{0.26\linewidth}p{0.20\linewidth}p{0.20\linewidth}p{0.20\linewidth}l@{}}
\toprule
Contrast & Crack IoU & Crack relF1@2 & Inactive IoU & Verdict \\
\midrule
MIM100-PF $-$ SUP100-PF (pool via encoder) & $+0.0048$ (0.0113) / $-0.0017$ (0.0100) & $+0.0035$ / $-0.0039$ & $-0.0118$ / $-0.0112$ & n.d. \\
SUP100-PF $-$ SUP100-g (recipe) & $-0.0012$ (0.0316) / $-0.0061$ (0.0100) & $+0.0117$ / $+0.0050$ & $+0.0150$ / $+0.0143$ & n.d. \\
CP100 $-$ SUP100-g, $\tau$ 0 (pool via paste) & $+0.0076$ (0.0307) / $+0.0003$ (0.0100) & $+0.0102$ / $+0.0012$ & $+0.0020$ / $+0.0007$ & n.d. \\
\bottomrule
\end{tabular}
\end{adjustbox}
\end{table}

\textbf{Verdict:} every contrast n.d., also with logit adjustment on both arms, so both pool routes and the recipe contrast were
closed at 100\% by the kill line. Beside it: PF kept the encoder (layer3 / layer4 cosine at 20,000 steps 0.928 / 0.936 against 0.015 /
0.001 for SUP100-g), so erasure does not explain the nulls; on E1 (descriptive, not blind) PF's crack recall was 0.213 against 0.169
and inactive IoU 0.245 against 0.041, all of the crack gain in the 29 scuff plus half-cut cells; CP100's decoy rate rose to 4.41
times the baseline's by 20,000 steps, so gate 6b's go rested on a 2,000 to 4,000-step fall that did not last. Our reading at the
time: the fine-tune recipe was the open lead, on E1 only.

\paragraph{Inference-only E1 checks (30 to 31 Oct).} A pre-registered read on validation and E1 (no test50) asked whether any
inference variant gained at least $+0.05$ E1 crack recall on the 29 scuff plus half-cut cells at a precision cost below 0.05. A five-seed
ResNet-50 ensemble passed narrowly ($+0.060$ at cost 0.047); test-time batch-norm re-estimation on half-cells gained $+0.174$ but cost
0.085 (fail); a ResNet-34 plus ResNet-50 ensemble gained $+0.036$ (fail); a 10-member ResNet-34 control gained $+0.001$, so the
ResNet-50 effect was not ensemble size.

\paragraph{Label work (30 to 31 Oct).} \emph{B1}: the five 50\%-label models audited the 239 train cells each never saw, flagging
components where all five were confident against the label; crack flags fell in 24 cells (10.0 per 100), not above the reviewed
validation base rate, so many flags were likely hard real cracks; three panel sheets went to the user, with no mask changed.
\emph{B2}: a ranked list of pool half-cells for the user to mark (scuff, crack, both, neither). Version 1 failed its own validation
(AUROC 0.328 and 0.396 on train scuff cells); version 2 (986 pool half-cells, ranked by predicted crack) was also unvalidated as a
scuff finder (AUROC 0.384 / 0.437) and went to the user as 90 ranked plus 10 blinded random panels. \emph{Large pool}: a public
unlabelled pool of 159,332 rows was checked and not downloaded: it has no E1 module family and no new half-cells, and its download
condition (a MIM100-PF gain) failed.

\section{A collaborator's eleven ideas and PF2 (31 Oct to 1 Nov)}\label{sec:midnov-ideas}

On 31 Oct (about 13:45) a collaborator accepted the protocol\_100 verdicts, approved PF2, and proposed eleven ideas aimed at unseen
panels, on the diagnosis that the remaining failure was domain generalisation rather than data quantity. The user asked for the
direction to be worked further without budget limits and with as little manual labelling as possible. Each idea was checked on CPU
by one agent and reviewed adversarially; only a reviewed pass could become a protocol. Table~\ref{tab:midnov-ideas} lists every idea
and its fate.

\begin{table}[htbp]
\centering
\scriptsize
\caption[A collaborator's eleven unseen-panel ideas (31 Oct to 1 Nov) and A1 to A6 (1 to 2 Nov)]{A collaborator's eleven
unseen-panel ideas (checked 31 Oct to 1 Nov) and six architecture ideas A1 to A6 (checked 2 Nov). CPU checks on validation, the
E-SCDD cells and earlier records; the checks of ideas 1 and 7 and the reviewer's checks also read E1 on 31 Oct, reads the ledger lists among those before
the E1 policy; A1 to A6 opened no test50 or E1 file. F51 / H36: the 51 full and 36 half-cut E-SCDD cells from modules not in
train (Section~\ref{sec:data-sets}).}
\label{tab:midnov-ideas}
\begin{tabular}{@{}p{0.25\linewidth}p{0.46\linewidth}p{0.21\linewidth}@{}}
\toprule
Idea & Main evidence & Outcome \\
\midrule
1 E1 threshold / calibration & E1 precision-recall curve below validation at every bias; label-free per-module thresholds lose on F51 and E1 & label-free thresholds dropped; half-cut routing: pre-check FAILED, parked \\
2 per-module intensity normalisation & E1 content brightness matches train (median 143 vs 138); the gap is the white frame of half-cut images; normalising lowers validation relF1@2 by 0.02 to 0.06 & dropped \\
3 neighbour-reference channel & 50\% of E1 cells have no same-module image anywhere & dropped; an in-crop variant parked \\
4 periodic-structure removal & the model already ignores busbars; crack vs scuff AUROC $\leq$ 0.45 & dropped alone; with 3 parked \\
5 synthetic half-cut cells & making validation cells half-cut drops FINAL's crack recall from 0.706 to 0.563 & passed (narrowed): run as HC \\
6 higher resolution & already run as Track M R5 (n.d.\ at 4$\times$ compute); E1 cracks are labelled wider, not thinner & dropped \\
7 cell-level and operating-point metrics & scorer validated; per-run bias unstable, so one pooled bias per arm & merged into every later protocol \\
8 foundation encoder & frozen probe: DINOv3 ConvNeXt-T F51 crack AP 0.363 against 0.310 for ImageNet ConvNeXt-T; nothing reached the 0.424 bar & passed to a 2-seed screen: A8 (KILL) \\
9 test-time adaptation & batch-norm adaptation gain mostly a calibration shift; per-panel variant T9 failed its $+0.03$ term and raised false alarms & dropped; test-time adaptation closed \\
10 AI-assisted labelling & E-SCDD crack labels: 75\% of components under 5 px; a segment-anything model fails on cracks; for inactive, SAM masks prompted by the model raised pooled IoU 0.4925 to 0.5653 & crack dropped; inactive passed as a method note, no arm \\
11 appearance augmentation & E1 brightness and contrast inside the house jitter range & no arm \\
\midrule
A1 two-branch encoder & needs idea 8, closed by A8; high-resolution branch is idea 6 & dropped \\
A2 instance norm in early layers & style premise contradicted (idea 2) & dropped \\
A3 geometry-aware input & HC already removes cut-line false alarms (123 to 3 px per run on H36) & dropped \\
A4 structured crack head & fragments are at most 5 to 7\% of non-E1 components & parked behind SKR2; not revived \\
A5 long-range context & no support & dropped \\
A6 cell gate plus pixel head & SKR's extra pixels are 99.6\% inside crack cells, so a gate cannot recover its precision & dropped as an arm \\
\bottomrule
\end{tabular}
\end{table}

Two data facts came out of these checks: E1 crack labels are about 1.5 times wider than train labels (3.66 against 2.41 px of area
per skeleton pixel), so E1 pixel recall partly reads width, and later protocols reported width-robust E1 metrics; and the crack
labels of the 36 E-SCDD half-cut cells are unusable (48\% of 1,792 px within 10\,px of a spacing, border or cut line), so no labelled
half-cut crack set exists outside E1.

\paragraph{PF2 (E1 opening \#1).} SUP100-g, SUP100-PF and RN50-PF (ResNet-50 under PF) on new seeds 5 to 9, second machine; trained 31 Oct
14:29 to 1 Nov 06:32, opened 06:33 to 08:26. Five co-primaries for PF $-$ house under rule 1: three on test50 and, pre-registered,
E1 crack recall and E1 inactive IoU. \textbf{Verdict: PASS.} E1 inactive IoU 0.0481 to 0.2371, $+0.1890$ (thr 0.0433) BETTER; E1
crack recall $+0.0271$ (thr 0.0370) n.d.; test50 relF1@2 $+0.0117$ / $+0.0120$ (thr 0.0100) BETTER in both sets; crack and inactive
IoU n.d. RN50-PF $-$ SUP100-PF: test50 inactive $-0.0289$ (thr 0.0251) WORSE in both sets: \textbf{FAIL}, ResNet-50 closed under PF.
Beside the verdict: the test50 pass clears the floor by 0.0017 and 0.0020, its bootstrap interval is $[-0.0034, +0.0292]$ and
leaving out the top 3 cells leaves $+0.0015$; all five PF seeds score 0.207 to 0.279 on E1 inactive against 0.025 to 0.089, with the
gain concentrated on half-cut cells (0.162 to 0.537). Later (1 Nov), the top 3 of the 7 source images of E1 inactive were found to
hold 87\% of its 9,158 pixels; without them PF2's gain is $+0.023$.

\paragraph{PF2E (E1 opening \#4, 1 Nov 19:02 to 19:13).} The deployable question: PF and house as 5-member mean-softmax ensembles
(seeds 5 to 9), inference only. The validation rule picked crack biases of 1.25 and 1.5, both of which dropped validation crack
precision by more than 0.05, so the verdict was read at bias 0 (argmax). \textbf{Verdict: PASS on E1 inactive IoU} (0.0665 to 0.2807,
$+0.2141$ $[+0.0303, +0.3820]$); E1 crack relF1@2 $+0.0355$ $[-0.0232, +0.0912]$ n.d.; test50 crack IoU, relF1@2 and inactive IoU
n.d.\ ($+0.0027$, $+0.0121$, $+0.0237$). Not independent of \#1 (same models, same cells). Beside it: half the gain is recall (0.304
to 0.627) and half fewer false inactive pixels on module-border and frame-edge pixels of three consecutive CSIR images; under a Bonferroni correction
the 99\% lower bound is about 0.000. The PF ensemble was recommended as the deployable candidate; its adoption was still open on
3 Nov.

\section{THIN-B3, A8, HC and B6 (1 to 2 Nov)}\label{sec:midnov-thin}

\paragraph{THIN-B3 (E1 opening \#2).} Three arms on SUP100-PF, seeds 5 to 9, against PF2's SUP100-PF runs (bit-identity shown over
2,000 steps): SKR, Skeleton Recall \citep{kirchhoff2024skeletonrecall} at weight 1; CLD, soft clDice \citep{shit2021cldice}
at weight 1; B3, scuff as a class of its own (the 649 scuff-as-crack pixels of Section~\ref{sec:midnov-escdd} become scuff).
Co-primaries: the three test50 primaries, for SKR and CLD also two operating-point metrics (component recall at component precision
$\geq$ 0.90 and relR@2 at relP@2 $\geq$ 0.90, at one validation bias per arm), and E1 crack relF1@2. Trained 1 Nov 06:32 to 19:39,
opened 19:40 to 21:13. \textbf{Verdicts:} SKR \textbf{FAIL} (test50 crack IoU $-0.0136$ / $-0.0140$, thr 0.0100, WORSE; component
recall $+0.0413$ BETTER; E1 relF1@2 $+0.0590$, thr 0.0392, BETTER); CLD \textbf{FAIL} (relR@2 at relP@2 $\geq$ 0.90 $-0.0374$ /
$-0.0376$, thr 0.015, WORSE); B3 \textbf{n.d.}\ on every co-primary, and the scuff class was never learned (train scuff IoU 0.000 in
all 5 seeds). Beside: SKR's predictions were about 1.5 times wider (test predicted to labelled pixels 0.95 to 1.54) and raised crack
on crack-free E1 cells $\times$2.4; its test50 WORSE is narrow (bootstrap $[-0.036, +0.012]$).

\paragraph{A8 screen (1 Nov 21:15 to 23:02).} Idea 8's kill gate: DINOv3 ConvNeXt-T \citep{simeoni2025dinov3} against ImageNet
ConvNeXt-T, both under PF, 2 seeds, 4,000 steps, scored on F51 and validation only (a data root whose test list was the validation
list kept test50 unread). GO needed F51 crack AP to gain with a 36-module bootstrap lower bound above 0. \textbf{Verdict: KILL}: $-0.0243$
(0.5448 against 0.5691), lower bound $-0.0463$; the foundation-encoder line was closed for the thesis. Beside: F51's crack is 76\% one
cell, and the ImageNet arm's own seed spread (0.067) exceeds the arm gap; KILL reads ``no sign of a gain'', not ``worse''.

\paragraph{HC (E1 opening \#3).} Synthetic half-cut views (cut line and gap drawn into labelled full cells) at probability 0.35 per
eligible sample (about 26\% of views), on SUP100-PF, seeds 5 to 9; trained 1 Nov 21:41 to 2 Nov 02:17, opened 02:18 to 03:10.
\textbf{Verdict: n.d.} on all five co-primaries: test50 crack IoU $+0.0092$ / $+0.0089$ (thr 0.0100), relF1@2 $+0.0026$ / $+0.0023$,
inactive $-0.0228$ / $-0.0230$ (thr 0.0289); E1 crack relF1@2 $+0.0135$ (thr 0.0429), E1 inactive $-0.0002$ (thr 0.0572). The
mechanism check passed (matched recall on synthetic half-cut validation views 0.8192 against 0.6959) and the F51 harm guard passed.
Beside: a seed-aware bootstrap put test50 inactive at $[-0.054, +0.005]$; our reading is that the simulator works on simulated
half-cuts while real E1 half-cuts did not move.

\paragraph{B6 (module-grouped cross-validation, 2 Nov).} Is PF's inactive gain also in distribution? PF against house, 5 folds over
train plus validation (532 cells) grouped by physical module, 2 seeds, 20 runs, trained on the training machine (all done 21:31) and
scored once on the second machine (22:11 to 22:13), by the user's decision on where each step could run. Primary: pooled out-of-fold inactive
IoU, pass if the two-stage bootstrap lower bound is above 0 and the gain at least 0.01. Power stated before the run: about 13\% at
$+0.017$. \textbf{Verdict: n.d.}, $-0.0144$ $[-0.0977, +0.0570]$. Beside: the realised standard error was 0.038 against a planned
0.020, so power at $+0.017$ was about 6\%; the interval does exclude an in-distribution gain as large as PF's E1 gain.

\section{SKR2, HC2, SKR2E and the B2 agent raters (2 Nov)}\label{sec:midnov-skr2}

\paragraph{SKR2 and HC2 (E1 opening \#5, 20:01 to 21:23).} SKR2 is Skeleton Recall at weight 0.1 on SUP100-PF; HC2 is HC with
inactive pixels set to ignore in synthetic views only. HC2's design went through four steps before lock (half dose; excluding
inactive-bearing sources; a pre-set retention rule; measured crack-pixel retention 0.399 below the 0.60 line, so the ignore
variant), and the user then chose to run it as built. Guards: G1, test50 inactive $\leq -0.015$ in both sets; G2 (SKR2), test50 relP@2 below $-0.03$ or E1 crack on the 74
crack-free cells above 1.5 times SUP100-PF; G3 (HC2), F51 crack AP upper bound below 0. \textbf{Verdicts:} SKR2 \textbf{PASS}: E1
crack relF1@2 $+0.0364$ (thr 0.0351) BETTER, the other six co-primaries n.d., no guard tripped; HC2 \textbf{FAIL} by G1 (test50
inactive $-0.0256$ / $-0.0259$), and its mechanism check failed (relP@2 0.90 never reached on synthetic validation views). Beside:
SKR2's margin was 0.0013 and the label survived 5 of 10 leave-one-seed-out drops, while a seed-aware interval of $[+0.004, +0.076]$
kept the effect; with 7 co-primaries the family-wise chance of a false BETTER is about 10\%. The route-closing rule for HC was
formally triggered; the user kept the half-cut route open and closed only the ignore design.

\paragraph{SKR2E (E1 opening \#6, 22:24 to 22:36).} The SKR2 and PF 5-member ensembles, verdict at a shared bias 0, with every guard
also checked at the headline pair (each ensemble at its own validation bias). The PF ensemble's headline bias was 0 (its width flag
tripped at 0.0753); SKR2's was 1.0 (validation precision drop 0.0485, under 0.05). At bias 0 all primaries passed or were n.d.\ (E1
relF1@2 $+0.0353$, test50 crack IoU $+0.0102$ $[+0.0031, +0.0197]$). At the headline pair G2a (test50 relP@2 $-0.0379$) and G2b (E1
crack on crack-free cells 5,358 against 2,528 px, $\times$2.12) tripped. \textbf{Verdict: FAIL}; the PF ensemble at argmax stayed the
recommended candidate. Beside: the width flag decided this by 0.0015; on test50, the PF ensemble at bias 1 equals the SKR2 ensemble
at bias 0, and \#6 re-measures \#5's models, so its bias-0 passes are no confirmation.

\paragraph{B2 agent raters.} The user asked why an agent could not mark the B2 panels. Two runs of one AI model were calibrated
blind before any B2 mark, under rules fixed in advance: name-free bundles, the key held by a reviewer, PASS/FAIL returned only; gate
A crack on 40 house-labelled CSIR train cells (sensitivity and specificity $\geq$ 0.70), gate B scuff on 32 pool half-cells of
E-SCDD (sensitivity $\geq$ 0.50, specificity $\geq$ 0.70). Both runs marked identically: crack gate PASS (20/20, 17/20), scuff gate
FAIL (sensitivity 1/8, needed 4/8). \textbf{No agent B2 marks were made}; B2 went back to the user, unmarked (state on 3 Nov).
Two runs of one model measure consistency, not independent reliability.

\section{SKR2C: trained 3 Nov, opened 14 Nov}\label{sec:midnov-skr2c}

SKR2C asked whether SKR2's E1 gain holds on fresh seeds 10 to 14 (10 runs on the second machine): a replication of seeds, not of
cells, and a third look at the same test50 and E1 cells for this contrast. The protocol was locked at 02:37 on 3 Nov. A code review
then found the unattended memory guard and queue not ready (seven must-fixes, among them a missing hold on outside GPU processes and a
missing re-check before a forced stop); after fixes and dry runs on dummies the run was launched at 03:06. Training read no test
cell (a data root with the test list set to validation). All 10 runs ended with return code 0 at 11:24 on 3 Nov; the guard never
fired. The record of 12 Nov notes that the laptop that drove the agents had been reinstalled and no agent survived. On 12 Nov
the user gave the OK to open
(``Open it now''); the opening set was frozen at 00:13 on 14 Nov and opened 00:22 to 03:54.

Two verdicts were registered. \textbf{R} (singles, rule 1, SKR2's seven co-primaries and guards G1, G2): \textbf{FAIL} by G1 (test50
inactive $-0.0198$ / $-0.0207$) and G2 (test50 relP@2 $-0.0780$ on all 50); E1 crack relF1@2 $+0.0289$ (thr 0.0281) BETTER, the other
six n.d. \textbf{D} (ensembles at a shared bias 0): \textbf{FAIL} by the inactive guard alone (test50 inactive $-0.0269$ / $-0.0283$);
E1 crack relF1@2 $+0.0342$ $[+0.0115, +0.0544]$ passes, E1 any-hit component recall $+0.0338$ $[0.0000, +0.0586]$ does not. The named
claim REPLICATED does not hold: its seed-aware interval was $[-0.0015, +0.0590]$ and R and D both failed. Fixed wording: ``the
fresh-seed ensemble pair differs on E1 relF1@2 (fixed-model CI; seed noise not included); replication at recipe level not shown''.
The SUP100-PF ensemble stayed, and the Skeleton Recall route closed for the deployable at weight 0.1. Beside: the G2 trip is one run
(seed 13) predicting 49,713 crack pixels on the noisy cell, where the label has 787; on clean 49 the relP@2 change is $-0.0119$. The
registered null false-FAIL chance of D's inactive guard alone was about 8\%; our reading, as in the verify record, is that a false
FAIL and a small real inactive cost cannot be told apart here.

\section{The E1 openings ledger, \#1 to \#7}\label{sec:midnov-ledger}

Table~\ref{tab:midnov-ledger} gives each registered E1 opening of the period with its rule and verdict, from the openings ledger
kept since 31 Oct; Table~\ref{tab:T17} has the short form and opening \#8. All seven read the same 106 cells, and their co-primaries
are uncorrected within and across openings. \#4 is not independent of \#1, \#6 not of \#5, and \#7 is the third look at the SKR2
contrast. PF2's opened SUP100-PF files were reused as comparator in \#2, \#3 and \#5, which is not a new read.

\begin{table}[htbp]
\centering
\scriptsize
\caption[E1 openings \#1 to \#7 with rule and verdict]{E1 openings \#1 to \#7 (second machine; same 106 E1 cells; test50 values
in-distribution, known leakage, context only). Verdict words as registered.}
\label{tab:midnov-ledger}
\begin{adjustbox}{max width=\linewidth}
\begin{tabular}{@{}p{0.03\linewidth}p{0.14\linewidth}p{0.13\linewidth}p{0.33\linewidth}p{0.29\linewidth}@{}}
\toprule
\# & Protocol & Opened & Rule (fixed before) & Verdict \\
\midrule
1 & PF2 & 1 Nov 06:33 to 08:26 & rule 1, 5 v 5; 3 test50 co-primaries (two sets) and E1 crack recall, E1 inactive IoU & PASS: E1 inactive $+0.1890$ BETTER; E1 crack recall n.d.; RN50-PF FAIL \\
2 & THIN-B3 & 1 Nov 19:40 to 21:13 & rule 1 per arm (a WORSE co-primary fails); test50, two operating-point metrics (SKR, CLD), E1 crack relF1@2 & SKR FAIL (E1 $+0.059$ BETTER); CLD FAIL; B3 n.d. \\
3 & HC & 2 Nov 02:18 to 03:10 & rule 1; test50 and E1 crack relF1@2, E1 inactive IoU; mechanism kill; F51 harm guard & n.d. (E1 $+0.0135$, $-0.0002$) \\
4 & PF2E & 1 Nov 19:02 to 19:13 & ensembles at headline bias; image / 71-cluster bootstrap; pass if interval above 0 & E1 inactive PASS ($+0.2141$); E1 crack relF1@2 n.d.; test50 n.d. \\
5 & SKR2\_HC2 & 2 Nov 20:01 to 21:23 & rule 1; SKR2 7 co-primaries with E1 relF1@2 and any-hit; guards G1, G2, G3 & SKR2 PASS ($+0.0364$, thr 0.0351); HC2 FAIL by G1 \\
6 & SKR2E & 2 Nov 22:24 to 22:36 & ensembles at shared bias 0; guards at bias 0 and the headline pair & FAIL by G2a and G2b at the headline pair \\
7 & SKR2C & 14 Nov 00:22 to 03:54 & R: rule 1 singles with guards; D: ensembles at bias 0 with guards; named claim needs a seed-aware interval above 0 and no FAIL & R FAIL; D FAIL; REPLICATED does not hold \\
\bottomrule
\end{tabular}
\end{adjustbox}
\end{table}

\paragraph{What the period led to.} Our reading: of the period's arms, only the PF recipe replicated an E1 gain on new seeds
(inactive IoU, on a 3-image readout), and the 5-seed ensemble gave the clearest test50 gain (context only); test50 and E1 had both been read too
often for any later choice to count as blind, and the 50-cell instrument could not resolve the crack effects in question. These are
the reasons the work moved, from 12 Nov, to the module-grouped benchmark of Chapter~\ref{ch:protocol}, with PF as its baseline and
one registered E1 opening for the final model (Chapter~\ref{ch:e1_result}).
